\section{Benchmark e Difesa
dell'Architettura}\label{benchmark-e-difesa-dellarchitettura}

Per concludere l'esposizione dell'architettura di BlueAgent, vengono
motivate le principali scelte infrastrutturali e di codice rimaste
aperte, difendendo la validità del progetto rispetto ai requisiti di un
Cyber Range orientato al patch management.

\begin{center}\rule{0.5\linewidth}{0.5pt}\end{center}

\subsection{Scelte di Piattaforma e
Infrastruttura}\label{scelte-di-piattaforma-e-infrastruttura}

\textbf{Perché Debian 13?} Offre un kernel moderno e stabile, garantendo
un'ottima compatibilita con i vari pacchetti utilizzati (Incus, OVN,
OVS, \ldots).

\textbf{Perche' ZFS?}: e' un file system generalmente ritenuto stabile e
affidabile che offre un sistema di CoW (Copy-on-Write) a basso costo,
senza il bisogno di duplicare i dati finche non vengono scritti. E'
stato scelto rispetto ad alternative come BTRFS poiche, ZFS e' ritenuto
più maturo e stabile.

\textbf{Perche' Incus?} A differenza di Docker, offre i container di
sistema, ossia container con un init system completo (e quindi un
proprio processo \passthrough{\lstinline!init!} con
\passthrough{\lstinline!PID!} 0). Permette ai concorrenti di operare
realisticamente, accedendo ai servizi di sistema ed utilizzare utility
complicate come \passthrough{\lstinline!apt!}, seppur in ambiente
containerizzato. Inoltre Incus supporta nativamente ZFS ed OVN.

\subsection{Difesa dell'Architettura}\label{difesa-dellarchitettura}

Il vantaggio piu grande che porta BlueAgent e' l'invisibilita totale
dell'Oracolo (o dello script di verifica), ribadiamo dunque che e'
impossibile per i partecipanti intercettare ed ascoltare traffico di
rete, o l'esecuzione dei comandi di verifica. La clonazione tramite CoW
in una sandbox OVN garantisce verifiche totalmente isolate, senza leak
di informazioni né disservizi sull'ambiente di produzione, con un uso di
risorse minimo.

La duplicazione tramite container rappresenta il giusto compromesso tra
l'uso di intere macchine virtuali e l'uso di namespace tramite NSJail
come visto in kCTF.

\begin{center}\rule{0.5\linewidth}{0.5pt}\end{center}

\subsection{Critica e Limiti
Architetturali}\label{critica-e-limiti-architetturali}

Nonostante i vantaggi, l'architettura presenta compromessi e vincoli
strutturali.

Combinare insieme il tunnel Geneve

\begin{itemize}
\tightlist
\item
  \textbf{UDP Path MTU Blackhole:} La combinazione di tunnel WireGuard
  (\(\text{MTU } 1420\)) e Geneve OVN (\(\text{MTU } 1442\), all'interno
  della \textbf{rete di verifica}) impone riduzioni di MTU. Mentre il
  TCP beneficia del MSS clamping via \passthrough{\lstinline!nftables!},
  i pacchetti UDP di grandi dimensioni con flag \emph{Don't Fragment},
  potrebbero essere scartati dal kernel.
\item
  \textbf{Limite scalabilità a 254 team:} L'indirizzamento
  deterministico \passthrough{\lstinline!10.k.x.x!} vincola \(k\) a un
  singolo byte (\(0 \le k \le 254\)), quindi ad un massimo di
  \(\text{255}\) team.
\item
  \textbf{Concorrenza e blocco multi-nodo:} Il meccanismo
  \passthrough{\lstinline!TwoTierLock!}
  (\passthrough{\lstinline!asyncio.Lock!} combinato con
  \passthrough{\lstinline!flock!}) è confinato al singolo nodo. Per
  scalare l'architettura su più nodi bisogna introdurre l'uso di un
  gestore di lock distribuito come Redis.
\end{itemize}
